% =====================================================================
% Note 8, hyperparameters. Moved out of the Supplementary Methods, 24 November (guide's
% comment: no Supplementary Methods; NMI allows algorithm detail as a Note).
% =====================================================================

\section{Every setting of the reported model}
\label{note:hyperparameters}

Table~\ref{stab:hyperparameters} lists every setting of the model reported in
the main text. Two of them, the information cost $w_{\mathrm{kl}}$ and the
adversarial weight $\lambda_{\mathrm{dab}}$, were varied systematically, and
their effects on the frontier are the subject of the main text and
\ref{note:frontier}. The rest were set during development, and no automated
hyperparameter search was used. In this note we give the reason for each
setting that was chosen against a measurement.

\textbf{Input.} On training steps, each cell's
measured genes were hidden from the encoder at random until 12{,}500
remained. This number lies just under the smallest panel in the corpus,
13{,}144 genes, so every cell shows the encoder the same number of
measured genes and a study cannot be recognised from how many genes it
measured. The decoder does not predict unmeasured genes, which are set
through the DNS sentinel. The 
token width and the 44 groups follow from the sizing argument of \ref{note:grouping}.

\textbf{Covariates are never withheld.} The run's recorded configuration
carries a probability of 0.25 for withholding the decoder's covariates on a
training step. That setting takes effect only together with a latent
component that no configuration in this paper uses, and the training metrics
confirm that no step withheld them. The table therefore gives the value in
force.

\textbf{Objective weights.} Each term other than reconstruction was
weighted so that its gradient on the latent stays below that of
reconstruction. We measured these gradients after the first epoch and
adjusted the weights against them, which reduced the menstrual-phase
classifier's weight from 500 to 40 and set the cell-similarity weight.

The adversary's capacity was set by its failure mode. An earlier
single-layer adversary of 64 units was beaten by the encoder without the
studies becoming inseparable. Its training loss fell to 0.21, against a
chance level of $\ln 11 = 2.40$, while the cells' nearest neighbours in the
latent still came from their own study 88\% of the time after balancing
for study size. The encoder had learned to fool that adversary, not to
remove study.

\textbf{Why the information-cost ramp is short.} By the end of the first
epoch the latent already held about 44 nats per cell in development. A ramp
that ends later prices information only after it has been stored, and a
cost cannot then remove it. The free-bits allowance never constrained the
latent, which stayed far above it throughout development.

\textbf{Learning rate and checkpoint selection.} The learning rate was
raised from an initial $3 \times 10^{-4}$ during development. The
checkpoint is selected on reconstruction rather than on a biological label
because the validation donors are uninformative about menstrual phase:
chance agreement on phase is 0.664 on the validation split, against 0.380
on the validation and test splits together. Every metric in the paper is
computed on the one selected checkpoint, and every other configuration
differs from it only in the settings named where it is introduced
(\ref{note:frontier} and~\ref{note:ablations}).

\textbf{Data.} Donors that one source deposited under two identifiers were
grouped before splitting (main Methods, \emph{Data collection}). Two deposits were excluded from training and
evaluation: one of chorionic villus and decidua, and one of organoids.

Overall, every setting in the table is the value in force during the
reported run, read from its recorded configuration and log, and every
other configuration in the paper differs from it only where stated.

\begin{table}[t]
\centering
\caption[Hyperparameters]{\textbf{Every setting of the reported model.}
Generated from the training run's own recorded configuration and log. KL,
Kullback--Leibler divergence. The adversary's gradient is reversed and
scaled by $\lambda_{\mathrm{dab}}$ on its way to the encoder.}
\label{stab:hyperparameters}
\small
% GENERATED by figures/scripts/supp_methods.py from run_config.json. Do not edit.
\begin{tabular}{@{}ll@{}}
\toprule
Setting & Value \\
\midrule
\multicolumn{2}{@{}l}{\textit{Input}} \\
\quad Count scaling & log, library size scaled to 10{,}000 \\
\quad Value encoding & one learned direction per gene, scaled by its count \\
\quad Token width $d$ & 32 \\
\quad Genes, groups & 15{,}496 genes in 44 groups \\
\addlinespace
\multicolumn{2}{@{}l}{\textit{Gene encoder}} \\
\quad Transformer layers per group & 2 \\
\quad Attention heads & 4 \\
\quad Dropout & 0.1 \\
\quad Summary tokens per group & 9 (1 learned, 8 pooled statistics) \\
\quad Size buckets & 4 \\
\addlinespace
\multicolumn{2}{@{}l}{\textit{Cell encoder}} \\
\quad Transformer layers & 2 \\
\quad Attention heads & 4 \\
\quad Dropout & 0.1 \\
\quad Width $d_{\mathrm{agg}}$ & 64 \\
\addlinespace
\multicolumn{2}{@{}l}{\textit{Latent}} \\
\quad Gaussian dimensions & 16 \\
\quad Circular dimensions & 0 \\
\addlinespace
\multicolumn{2}{@{}l}{\textit{Decoder}} \\
\quad Width $d_{\mathrm{dec}}$ & 64 \\
\quad Tissue, platform, study embeddings & 16, 16, 16 dimensions \\
\quad Covariates zeroed per step & probability 0.25 \\
\addlinespace
\multicolumn{2}{@{}l}{\textit{Objective}} \\
\quad Reconstruction weight $w_{\mathrm{nb}}$ & 1 \\
\quad KL weight $w_{\mathrm{kl}}$ & 1 \\
\quad KL ramp & 0 to 0.5 over the first 5\% of steps \\
\quad KL free bits & 0.25 \\
\quad Adversary weight $w_{\mathrm{dab}}$, reversal $\lambda_{\mathrm{dab}}$ & 0.5, 168 \\
\quad Adversary network & 3 layers, 256 hidden units \\
\quad Phase head weight, label smoothing & 40, 0.05 \\
\quad Cell-similarity weight, threshold & 50, 0.8 \\
\addlinespace
\multicolumn{2}{@{}l}{\textit{Optimisation}} \\
\quad Optimiser & AdamW, $\beta = (0.9, 0.999)$ \\
\quad Learning rate, weight decay & $10^{-3}$, $10^{-4}$ \\
\quad Schedule & cosine over 20 epochs \\
\quad Batch size & 256 \\
\quad Gradient clipping & norm 1 \\
\quad Precision & bf16 \\
\quad Early stopping & held-out reconstruction, patience 5 after 2 epochs \\
\quad Selected epoch & 4 (held-out reconstruction 0.4024) \\
\addlinespace
\multicolumn{2}{@{}l}{\textit{Data}} \\
\quad Split & by donor, 70/15/15, seed 24 \\
\quad Training seed & 24 \\
\quad Parameters & 1{,}883{,}986 \\
\bottomrule
\end{tabular}

\end{table}
